\chapter*{Lời Mở Đầu}
\addcontentsline{toc}{chapter}{Lời Mở Đầu}
\markboth{Lời Mở Đầu}{Lời Mở Đầu}

\begin{truyen}{Lời Mở Đầu}{Opening Words}
\chuhoa{C}{ó những cuốn sách viết để kể chuyện cho người khác nghe. Cuốn này được viết trước hết để nhắc chính người viết phải tỉnh táo hơn mỗi lần cầm đồng tiền trong tay, bởi người viết đã từng lớn lên trong một tuổi thơ không có quyền lãng phí.}
\textit{(Some books are written to tell stories to others. This one is written first to remind the writer himself to stay awake each time he holds money in his hand, for he grew up in a childhood that had no right to waste.)}

Năm bảy tuổi, người viết mất cha. Gia đình không có nhà cửa ổn định, lưu lạc nhiều nơi, từng có những ngày phải đi chăn vịt, đi học muộn hơn bạn bè, sớm hiểu thế nào là tủi thân và bấp bênh. Chính những năm tháng ấy đáng lẽ phải làm cho con người biết quý từng đồng tiền suốt đời. Nhưng lạ thay, ngay cả người từng khổ vẫn có thể quên nhanh cái giá của đồng tiền khi đời sống bận rộn, trách nhiệm chồng chất và mệt mỏi bào mòn sự tỉnh táo.
\textit{(At the age of seven, the writer lost his father. The family had no stable home, wandered from place to place, lived through days of herding ducks, entered school later than others, and learned early what humiliation and uncertainty felt like. Those years should make a person value every coin for life. Yet strangely, even those who once suffered can quickly forget the cost of money when life grows busy, burdens multiply, and exhaustion erodes clarity.)}

Vì thế, cuốn sách này không tô vẽ sự giàu có. Nó chỉ cố giữ lại vài nguyên tắc giản dị: biết phân biệt điều cần và điều muốn, biết để dành cho ngày biến cố, biết sống vừa sức, biết dạy con bằng nếp sống của mình, biết dùng đồng lương giáo viên và những giờ dạy thêm đúng chỗ, và biết dừng lại trước những ham muốn chỉ làm lòng hả một lúc mà ví tiền đau rất lâu.
\textit{(For that reason, this book does not glorify wealth. It simply tries to preserve a few plain principles: distinguishing between need and desire, saving for difficult days, living within one’s means, teaching children through example, using a teacher’s salary and extra tutoring income properly, and stopping before cravings that satisfy the heart briefly but hurt the wallet for a long time.)}

Mười chương sau đây được viết theo giọng truyện kể gần với hồi ký hơn, xen điển tích cũ để soi vào việc sống hôm nay. Ở cuối mỗi chương, người viết chủ ý thêm hai lớp bài học: một lớp để tự răn chính mình, và một lớp để có thể mang vào lớp học, giúp học sinh hiểu rằng đạo đức không nằm ngoài chuyện dùng tiền, biết ơn và tự giữ mình. Nếu đọc xong mà chỉ giữ lại được một thói quen nhỏ tốt hơn, thì cuốn sách này đã không vô ích.
\textit{(The ten chapters that follow are written in a narrative voice closer to memoir, woven with old anecdotes to illuminate present-day living. At the end of each chapter, two layers of lessons are added deliberately: one for self-discipline, and one that can be carried into the classroom so students may understand that morality is not separate from how one uses money, gratitude, and self-restraint. If, after reading, even one small better habit remains, then this book has not been useless.)}
\end{truyen}

\begin{baihoc}
Tiết kiệm không phải là câu chuyện của tiền trước hết. Nó là câu chuyện của trí nhớ, của kỷ luật, và của lòng biết sợ những ngày gia đình mình có thể chông chênh. Với một người thầy, nó còn là cách dạy học sinh bằng chính đời sống của mình.
\textit{(Saving is not primarily a story about money. It is a story about memory, discipline, and the healthy fear of days when one’s family may become unstable. For a teacher, it is also a way of teaching students through one’s own life.)}
\end{baihoc}

\begin{ghinhoanh}
\textbf{Short English to remember:}

Saving begins in memory.

Discipline keeps tomorrow from collapsing.
\end{ghinhoanh}

\begin{tuhocnhanh}
1. Điều gì trong đời khiến em quý đồng tiền hơn trước? \textit{What in life has made you value money more than before?}

2. Em đang tiêu tiền theo thói quen nào mà chưa dám nhìn thẳng? \textit{What spending habit are you still unwilling to face honestly?}

3. Sau khi đọc quyển này, em muốn đổi điều gì đầu tiên? \textit{After reading this book, what is the first thing you want to change?}
\end{tuhocnhanh}

\ngancach